%%%%%%%%%%%%%%%%%%%%%%%%%%%%%%%%%%%%%%%%%%%%%%%%%%%%%%%%%%%%%
% ==================== 第三章：模型假设 ====================
% BZD 数模社制作——国赛 B 题通用模板
% 使用时请根据题目条件和后续模型逐一替换【】中的内容；假设数量可增减。
%
% ✓ 自查清单：
%   □ 是否只保留后续建模真正需要、且题目没有直接给出的假设？
%   □ 每条假设是否说明依据、简化对象、适用范围及其对模型的作用？
%   □ 是否保留核心机理、物理约束和业务约束，只忽略次要因素？
%   □ 数据假设是否与附件质量检查和实际预处理方法一致？
%   □ 是否避免“数据完全真实可靠”“误差均可忽略”等无依据断言？
%   □ 每条假设是否在后续公式、参数、约束或求解过程中实际使用？
%   □ 若假设不成立，是否能够说明其可能影响和替代处理方式？
%   □ 是否删除重复假设和仅为计算方便而设置的不合理假设？
%%%%%%%%%%%%%%%%%%%%%%%%%%%%%%%%%%%%%%%%%%%%%%%%%%%%%%%%%%%%%

\section{模型假设}

在保留问题核心机理和主要约束的前提下，为明确研究边界并使模型具有可求解性，本文根据题目条件和各问建模需要作出如下必要假设：

\begin{enumerate}
  \setlength{\itemsep}{0.8em}

  \item 假设【具体内容】。

  \textbf{依据与作用：}【说明该假设的题设依据、数据依据、客观规律或合理来源；指出被简化的对象、适用的问题及其对变量、参数、约束或求解过程的作用。】

  \item 假设【具体内容】。

  \textbf{依据与作用：}【说明该假设的题设依据、数据依据、客观规律或合理来源；指出被简化的对象、适用的问题及其对变量、参数、约束或求解过程的作用。】

  \item 假设【具体内容】。

  \textbf{依据与作用：}【说明该假设的题设依据、数据依据、客观规律或合理来源；指出被简化的对象、适用的问题及其对变量、参数、约束或求解过程的作用。】

  \item 假设【具体内容】。

  \textbf{依据与作用：}【说明该假设的题设依据、数据依据、客观规律或合理来源；指出被简化的对象、适用的问题及其对变量、参数、约束或求解过程的作用。】

  \item 假设【具体内容】。

  \textbf{依据与作用：}【说明该假设的题设依据、数据依据、客观规律或合理来源；指出被简化的对象、适用的问题及其对变量、参数、约束或求解过程的作用。】
\end{enumerate}
